\documentclass[12pt, a4paper, oneside]{book}
%\author{aya}
\usepackage[top=3cm,right=2cm,bottom=3cm,left=2cm]{geometry}
\usepackage{graphicx}
\usepackage{indentfirst}
\usepackage{anyfontsize}
\usepackage{wrapfig}
\usepackage{enumerate,letltxmacro}
\LetLtxMacro\itemold\item
\graphicspath{{images/}}

\setlength{\parindent}{4em}
\renewcommand{\item}{\itemindent4.5em\itemold}

\begin{document}
	\begin{center}
			\fontsize{16}{24}{
				\textbf{
					Pemrograman Berbasis Kerangka Kerja\\
					Sistem Pendukung Domain Kesehatan\\
		}
	}
	\end{center}
	\begin{center}
		\includegraphics[width=0.5\textwidth]{lambangits}
		\linebreak
			\begin{table}[h]
				\centering
				\begin{tabular}{lr}
					M. Vijay Fathur       & 5112100043 \\
					Nabila Tsurayya Silmi & 5112100055 \\
					Faishal Azka J        & 5112100061 \\
					R. M Iskandar Z       & 5112100101 \\
					Reyhan Arief          & 5112100175 \\
				\end{tabular}
			\end{table}
		\end{center}
		\begin{center}
		{		
			\fontsize{16}{24}{
				\textbf{
					Fakultas Teknologi Informasi\\
					Jurusan Teknik Informatika\\			
					Institut Teknologi Sepuluh Nopember Surabaya\\
				} 
			}
		}
	\end{center}
	
	\thispagestyle{empty}
	\pagebreak
	\tableofcontents
	\pagebreak
	\listoffigures
	\pagebreak
	\chapter{Deskripsi Produk}
	\section{Perspektif Produk}
		Produk yang yang hendak dikembangkan adalah sistem pendukung kesehatan, yaitu Sistem Asuransi Kesehatan.\par
		{Tujuan dibuatnya Sistem Asuransi Kesehatan antara lain, adalah sebagai berikut:
		\begin{enumerate}
			\item Mempermudah proses pendaftaran anggota asuransi kesehatan.
			\item Mempermudah administrasi asuransi kesehatan.
			\item Mempermudah anggota untuk mendapatkan informasi lokas fasilitas kesehatan.
		\end{enumerate}}
		\par
		Diagram Konteks
	\section{Klasifikasi dan Karakteristik Pengguna}
		\begin{table}[h]
		\centering
		\begin{tabular}{c | c | l | c}
			\hline
			No	& Kategori Pengguna	& Hak Akses ke Aplikasi 						& Kemampuan User\\
			\hline
			1.	& Calon Anggota		& Mendaftarkan diri sebagai anggota.			& -\\
			\hline
			2.	& Anggota			& 1. Melihat fasilitas kesehatan terdekat.		& -\\
				&					& 2. Melihat status iuran.						& \\
			\hline
			3. 	& Administrator		& 1. Melihat data anggota						& Mengerti lalala \\
		\end{tabular}
	\end{table}
	\section{Lingkungan Pengoperasian}
	
	\section{Batasan dan Desain Implementasi}
	\section{Asumsi dan Ketergantungan}
	\pagebreak
	\chapter{Fitur}	
	\end{document}
\end{document}
